\vspace{-20pt}

{
\vspace*{-0.2cm}
\centering
\includegraphics[width=\linewidth, trim=0 0 0 0, clip]{figures/teaser.pdf}
\vspace{-6pt}
\captionof{figure}{\sotopia: An \textbf{open-ended social interaction environment}. In each episode, \sotopia first samples a social scenario context, goals, and characters, and then assigns a social goal and character to each agent involved. Agents (artificial agents or humans) in \sotopia role-play characters while attempting to achieve their goals. The agents' performance is evaluated through a multi-dimensional framework, \sotopiaeval.}
\label{fig:teaser}
\vspace{10pt}
}

\begin{abstract}
\vspace{-10pt}
\emph{Humans are social beings}; we pursue social goals in our daily interactions, which is a crucial aspect of social intelligence. Yet, AI systems' abilities in this realm remain elusive.
We present \sotopia, an open-ended environment to simulate complex social interactions between artificial agents and evaluate their social intelligence.
In our environment, agents role-play and \textit{interact} under a wide variety of scenarios;
they coordinate, collaborate, exchange, and compete with each other to achieve complex social goals. 
We simulate the role-play interaction between LLM-based agents and humans within this task space and evaluate their performance with a holistic evaluation framework called \sotopiaeval.
With \sotopia, we find significant differences between these models in terms of their social intelligence, and we identify a subset of \sotopia scenarios, \sotopia-hard, that is generally challenging for all models. 
We find that on this subset, GPT-4 achieves a significantly lower goal completion rate than humans and struggles to exhibit social commonsense reasoning and strategic communication skills. 
These findings demonstrate \sotopia's promise as a general platform for research on evaluating and improving social intelligence in artificial agents. 
\end{abstract}


